\section{Verdict}

The daylight reading survives as a serious, conditional reading of the chapter.
If Genesis 15:1--18 narrates one continuous episode in its presented order,
God's command to look at the stars precedes the sunset announced at verse 12.
The Hebrew and Vulgate support that sequence; no explicit interval interrupts
it.

The condition matters. Vision can loosen the relation between narrated and
waking time. Biblical narrative can omit an interval. Swete's identified
printed witness reduces the chapter's temporal closure. The verified
historical witnesses preserve more than one resolution, so the answer cannot
be treated as self-evident.

The strongest conclusion is textual, not astronomical: Genesis presents the
star promise before its own sunset. Whether the represented world is daylight,
visionary time, a second night, supernatural sight, or cosmological transport
is interpretation. The daylight-belief meditation belongs to this project's
synthesis and remains explicitly conditional.

